\booksection{第一章附：鴻蒙作業系統專題}

\chapterintro{鴻蒙（HarmonyOS）係華為自主研發嘅作業系統，由2019年發佈至今，已經走過咗幾個大版本。佢唔只係「又一個安卓替代品」，而係有住獨特設計理念嘅新物種。等我哋一齊深入瞭解呢個正在崛起嘅作業系統。}

\booksubsection{一、鴻蒙係乜嘢？點解佢與眾不同？}

鴻蒙嘅核心理念係「萬物互聯」。唔同於傳統作業系統主要服務於單一設備，鴻蒙由設計之初就考慮咗多設備協同。

\textbf{分佈式架構}係鴻蒙嘅靈魂。簡單講，佢可以將多個設備虛擬成一個「超級終端」。你嘅手機、平板、電腦、手錶、電視，唔再係孤立嘅個體，而係一個整體嘅「器官」，可以互相調用對方嘅能力。

舉個例：你喺度用手機睇影片，想喺大螢幕上睇，傳統方式係投屏；而喺鴻蒙上，你可以直接將影片「流轉」到電視上繼續睇，手機變成遙控器。呢個唔係簡單嘅投屏，而係應用本身喺度跨設備運行。

\booksubsection{二、鴻蒙4.3：基於AOSP嘅成熟版本}

鴻蒙4.3係目前華為手機上廣泛部署嘅版本。需要說明嘅係，鴻蒙4.3兼容安卓開源項目（AOSP）嘅代碼框架，因此佢可以完整兼容安卓應用——你喺華為應用市場下載嘅大部分App，本質上仲係安卓應用，通過鴻蒙嘅兼容層運行。但鴻蒙4.3並非簡單套殼安卓，佢喺AOSP基礎上加入咗大量華為自研嘅分佈式能力、方舟引擎等技術。

\subsubsection*{核心特性：}

\begin{itemize}
  \item \textbf{安卓應用兼容}：由於基於AOSP開發，鴻蒙4.3可以正常運行幾乎所有安卓應用，用戶幾乎感受唔到差異。
  \item \textbf{方舟引擎}：華為自研嘅編譯器同運行時，等應用運行更流暢、更慳電。相比早期版本，流暢度提升明顯。
  \item \textbf{超級終端}：喺控制中心，可以直觀咁睇到周圍嘅其他鴻蒙設備，拖拽即可連接。手機調用電腦嘅鍵盤滑鼠、平板調用手機嘅鏡頭，都好方便。
  \item \textbf{服務卡片}：類似iOS嘅小組件，但更靈活。可以上滑App圖標展開卡片，顯示關鍵資訊，唔使打開App。
  \item \textbf{純淨模式}：淨係允許由官方應用市場安裝應用，減少惡意軟件風險。
  \item \textbf{私隱中心}：集中管理應用權限，查看應用行為記錄。
\end{itemize}

\subsubsection*{用戶介面特點：}

\begin{itemize}
  \item \textbf{雜誌鎖屏}：鎖屏介面顯示精美圖片同文章，劃走即進入桌面。
  \item \textbf{萬能卡片}：桌面圖標可以展開成卡片，顯示更多資訊。
  \item \textbf{超級中轉站}：一個臨時存儲區，可以將文字、圖片、檔案暫存，然後拖到任意App中使用。跨應用分享變得無比絲滑。
  \item \textbf{智慧搜尋}：全局搜尋，可以搜應用、檔案、設定、網頁。
\end{itemize}

\booksubsection{三、鴻蒙NEXT（鴻蒙6）：華為自主作業系統嘅真正起點}

鴻蒙NEXT（HarmonyOS NEXT 5.0）同鴻蒙4.3有住根本性嘅區別——佢採用華為自研嘅鴻蒙微內核，唔再依賴Linux內核同AOSP代碼，都\textbf{唔再直接支援安卓應用}。呢個係鴻蒙由「兼容安卓嘅過渡形態」正式走向「完全獨立嘅作業系統」嘅關鍵一步。

\subsubsection*{重大變革：}

\begin{itemize}
  \item \textbf{完全自主內核}：鴻蒙NEXT使用華為自研嘅內核，唔再包含任何AOSP代碼。系統更精簡、更安全、更高效。
  \item \textbf{唔再支援安卓應用}：呢個係最大嘅變化。鴻蒙NEXT淨係可以運行原生鴻蒙應用，無法安裝同運行APK檔案。升級前請確認你常用嘅App已經有鴻蒙原生版本。
  \item \textbf{原生應用生態}：華為正在全力推動開發者將應用遷移到鴻蒙原生平臺。目前主流App如微信、支付寶、淘寶等已推出鴻蒙原生版本，但部分小眾應用可能尚未適配。
  \item \textbf{星盾安全架構}：全新嘅安全體系，由硬件到軟件嘅全鏈路保護。支援更細粒度嘅權限控制，比如「模糊定位」——應用淨係可以知道你大概喺邊，唔知道精確位置。
  \item \textbf{AI深度融合}：系統級AI能力，唔只係某個App有AI。比如AI字幕可以實時翻譯影片，AI摳圖可以智能識別主體。
  \item \textbf{跨設備無縫流轉}：更進一步嘅分佈式能力。應用狀態可以喺設備間無縫切換，你喺手機上編輯嘅文檔，行到電腦前自動喺電腦上打開。
  \item \textbf{元服務}：無需安裝即可使用嘅輕量級服務。比如掃二維碼直接打開某個功能，唔使先下載App。
\end{itemize}

\begin{tcbwarn}{升級鴻蒙NEXT前必讀}
升級到鴻蒙NEXT後，你嘅手機將無法運行安卓應用。請務必喺升級前確認：(1)你常用嘅App已有鴻蒙原生版本；(2)你嘅工作或生活中冇必須依賴安卓應用嘅場景。升級後無法回退到鴻蒙4.3。

\end{tcbwarn}

\subsubsection*{操作邏輯變化：}

\begin{itemize}
  \item \textbf{手勢導航}：左右邊緣滑動返回，底部上滑返桌面，上滑停頓入多任務。
  \item \textbf{智慧多窗}：支援分屏、懸浮窗、平行視界等多種多任務形態。
  \item \textbf{暢連通話}：華為嘅視像通話服務，支援螢幕共享、美顏、背景虛化。
\end{itemize}

\booksubsection{四、鴻蒙生態：設備與應用}

鴻蒙唔只係手機系統，而係一個覆蓋多種設備嘅生態系統。

\subsubsection*{支援嘅設備類型：}

\begin{itemize}
  \item \textbf{手機與平板}：核心設備，功能最完整。
  \item \textbf{智能手錶與手環}：獨立運行或同手機協同。
  \item \textbf{智慧屏（電視）}：家庭娛樂中心。
  \item \textbf{智能音箱}：語音互動入口。
  \item \textbf{汽車}：鴻蒙座艙，智能駕駛艙體驗。
  \item \textbf{智能家居}：雪櫃、冷氣、洗衣機等。
\end{itemize}

\subsubsection*{應用生態現狀：}

鴻蒙應用市場已有數百萬應用，主流App基本都有。華為都喺度積極推動開發者開發原生鴻蒙應用。截至2026年中，微信、支付寶、淘寶、抖音、美團、京東等頭部應用均已推出鴻蒙原生版本，日常使用基本無障礙。部分小眾應用同行業應用可能仲喺度適配中，升級前建議確認你常用嘅App係咪已支援。

\booksubsection{五、鴻蒙實用技巧}

\begin{itemize}
  \item \textbf{超級終端}：由螢幕右上角下滑打開控制中心，點擊「超級終端」，可以睇到周圍嘅其他設備，拖拽連接。
  \item \textbf{超級中轉站}：長按文字/圖片/檔案，選擇「複製到中轉站」，然後喺其他App中長按貼上，或拖拽出嚟。
  \item \textbf{服務卡片}：上滑App圖標展開卡片，長按卡片可以調整大小或添加到桌面。
  \item \textbf{智慧多窗}：由螢幕左右邊緣向內滑動並停頓，打開側邊欄，點擊應用以懸浮窗打開。
  \item \textbf{私隱保護}：設定 > 私隱，可以查看應用權限使用記錄，設定模糊定位等。
\end{itemize}

\begin{tcbwarn}{注意事項}
鴻蒙系統喺度不斷更新迭代，唔同版本可能有啲差異。本文描述基於鴻蒙4.3同6嘅主要特性，具體功能請以你嘅手機實際版本為準。建議保持系統更新，以獲得最新功能同安全補丁。

\end{tcbwarn}

\begin{tcbpractice}{而家就試：體驗鴻蒙特色功能}
\begin{enumerate}
  \item 打開控制中心，揾到「超級終端」，睇下週圍有冇其他鴻蒙設備。
  \item 試下「超級中轉站」：複製一段文字到中轉站，然後喺備忘錄中拖出嚟。
  \item 上滑幾個常用App圖標，睇下佢哋嘅服務卡片顯示乜嘢內容。
  \item 去設定 > 私隱，查看最近邊啲應用使用咗敏感權限。
\end{enumerate}

\end{tcbpractice}
